\chapter{TYPE INFERENCE ALGORITHM}

\section{Overview}
Type inference in this project means that the programmer does not write explicit
type declarations for variables or functions. Instead, the compiler infers
types from assignments, literals, operations, function calls, and return
statements. This keeps the language simple while still preserving static type
safety. Specifically, the algorithm is \emph{syntax-directed} and local: each
expression's type is determined by its subexpressions and operator token alone,
without unification variables or constraint solving.

The type checker runs after parsing. It receives the AST and validates the
program before execution. Any inconsistent use of types is reported as a type
error, and execution is stopped.

This implementation is not Hindley--Milner inference; it uses a
syntax-directed local strategy with first-call parameter fixation for called
functions and a body-usage post-pass for never-called functions.

The interpreter stage does not repeat static type validation: it walks the same AST
with ordinary Python representations of values. Static type information remains
in the symbol table produced by the checker; run-time values are not separately
labelled with \texttt{LanguageType}, except that the built-in \texttt{print}
formats output
by inspecting the value's Python type.

\section{Type Environment}
The type environment is represented by the symbol table. Conceptually, it is a
mapping

\begin{center}
$\Gamma : \textit{Identifier} \rightarrow \textit{Type}$
\end{center}

The environment grows as the checker visits the AST, and on initial binding a
variable name is inserted with its inferred type. When a new block or a
function call is entered, a child environment is created so that nested scopes
can be checked independently while still accessing outer bindings when needed.

\section{Type Inference Procedure}

The type checker performs type inference by traversing the abstract syntax tree
before interpretation. The algorithm is syntax-directed: each AST node is
visited, the types of its child nodes are inferred first, and the current node
is checked according to the language typing rules.

\begin{center}
\textbf{Algorithm} \textsc{TypeCheck}$(\textit{node},\; \Gamma)$
\end{center}

\noindent\textbf{Input:} \textit{node} = current AST node;\quad
$\Gamma$ = current type environment / symbol table

\noindent\textbf{Output:} inferred type of \textit{node}, or \texttt{TypeError}

\begin{enumerate}[leftmargin=1.5em, label=\arabic*., itemsep=2pt]
  \item \textbf{If} \textit{node} is a \textbf{Literal}:\quad
        return the corresponding primitive type:
        \texttt{Integer}, \texttt{Float}, \texttt{Boolean}, or \texttt{String}.

  \item \textbf{If} \textit{node} is an \textbf{Identifier}:\quad
        look up the identifier in $\Gamma$.
        If it is not found, report an undefined variable error.
        Otherwise, return its stored type.

  \item \textbf{If} \textit{node} is a \textbf{BinaryOp}:
        \begin{itemize}[nosep, leftmargin=1.5em]
          \item Infer the type of the left operand and the right operand.
          \item If the operator is \texttt{+}, \texttt{-}, \texttt{*}, or \texttt{/}:\quad
                require both operands to be \texttt{Integer}; return \texttt{Integer}.
          \item If the operator is \texttt{+.}, \texttt{-.}, \texttt{*.}, or \texttt{/.}:\quad
                require both operands to be \texttt{Float}; return \texttt{Float}.
          \item If the operator is \texttt{==} or \texttt{!=}:\quad
                require both operands to be the same arithmetic type; return \texttt{Boolean}.
          \item Otherwise, report a \texttt{TypeError}.
        \end{itemize}

  \item \textbf{If} \textit{node} is an \textbf{Assignment}:
        \begin{itemize}[nosep, leftmargin=1.5em]
          \item Infer the type of the right-hand-side expression.
          \item If the variable is not yet in $\Gamma$: insert it with the inferred type.
          \item Else: verify the new type matches the existing type; if not, report a \texttt{TypeError}.
        \end{itemize}

  \item \textbf{If} \textit{node} is an \textbf{If} statement:
        \begin{itemize}[nosep, leftmargin=1.5em]
          \item Infer the condition type; require it to be \texttt{Boolean}.
          \item Type-check the then-block in a child environment.
          \item If an else-block exists, type-check it in a child environment.
        \end{itemize}

  \item \textbf{If} \textit{node} is a \textbf{While} statement:
        \begin{itemize}[nosep, leftmargin=1.5em]
          \item Infer the condition type; require it to be \texttt{Boolean}.
          \item Type-check the loop body in a child environment.
        \end{itemize}

  \item \textbf{If} \textit{node} is a \textbf{FunctionDef}:
        \begin{itemize}[nosep, leftmargin=1.5em]
          \item Register the function name in $\Gamma$.
          \item Parameter types start as placeholders; for called functions they
                are fixed by the first call site, while never-called functions
                are assigned using body-usage inference in the post-pass.
          \item Return type is inferred from \texttt{return} statements in the body.
        \end{itemize}

  \item \textbf{If} \textit{node} is a \textbf{FunctionCall}:
        \begin{itemize}[nosep, leftmargin=1.5em]
          \item If the function is the built-in \texttt{print}:\quad
          require exactly one argument, infer that argument type,
          reject \texttt{Void} arguments, and return \texttt{Void}.
          \item Otherwise, infer all argument types, validate them against the
                function signature, type-check the callee body if needed, and
                return the inferred function return type.
        \end{itemize}
\end{enumerate}

\subsection{Example type inference trace}
Table~\ref{tab:type-inference-trace} traces the algorithm step by step for
the program fragment

\begin{lstlisting}
x = 3;
y = 4;
z = x + y;
\end{lstlisting}

\noindent showing how the type environment $\Gamma$ grows as each
sub-expression is inferred.

\begin{table}[H]
\centering
\caption{Step-by-step type inference trace for \texttt{x = 3; y = 4; z = x + y;}.\label{tab:type-inference-trace}}
\small
\begin{tabular}{clll}
\toprule
Step & Expression & Inferred type & Environment $\Gamma$ \\
\midrule
1 & \texttt{3}         & Integer & $\emptyset$ \\
2 & \texttt{x = 3}     & Integer & $\{x\!:\!\textit{Integer}\}$ \\
3 & \texttt{4}         & Integer & $\{x\!:\!\textit{Integer}\}$ \\
4 & \texttt{y = 4}     & Integer & $\{x\!:\!\textit{Integer},\ y\!:\!\textit{Integer}\}$ \\
5 & \texttt{x + y}     & Integer & unchanged \\
6 & \texttt{z = x + y} & Integer & $\{x\!:\!\textit{Integer},\ y\!:\!\textit{Integer},\ z\!:\!\textit{Integer}\}$ \\
\bottomrule
\end{tabular}
\end{table}

\section{Inference Rules}

\subsection{Judgement Forms}
The static semantics are written using the following judgement forms:
\[
\Gamma \vdash e : T
\]
which reads ``under type environment $\Gamma$, expression $e$ has type $T$'', and
\[
\Gamma \vdash e : \texttt{TypeError}
\]
which denotes a statically rejected expression.

\subsection{Literals}
Literal nodes have direct types:

\[
\frac{}{\Gamma \vdash 42 : \texttt{Integer}}\ \textsc{(T-IntLit)}
\qquad
\frac{}{\Gamma \vdash 3.14 : \texttt{Float}}\ \textsc{(T-FloatLit)}
\]
\[
\frac{}{\Gamma \vdash \texttt{true} : \texttt{Boolean}}\ \textsc{(T-BoolLit)}
\qquad
\frac{}{\Gamma \vdash \texttt{"hi"} : \texttt{String}}\ \textsc{(T-StringLit)}
\]

\subsection{No Operator Overloading}
The language does not support operator overloading. Following the Caml~Light
convention, integer and float arithmetic use entirely separate operator tokens,
so every operator has a single, fixed type signature:
\begin{itemize}
  \item \texttt{+}, \texttt{-}, \texttt{*}, \texttt{/} --- require two \texttt{Integer} operands; result is \texttt{Integer}.
  \item \texttt{+.}, \texttt{-.}, \texttt{*.}, \texttt{/.} --- require two \texttt{Float} operands; result is \texttt{Float}.
  \item \texttt{+} is not overloaded for string concatenation; applying it to a \texttt{String} is a type error.
  \item \texttt{==} and \texttt{!=} are not overloaded for \texttt{Boolean} or \texttt{String} equality.
  \item No implicit numeric promotion: \texttt{Integer + Float} is a type error.
\end{itemize}
As every operator has a unique, unambiguous type derivable from its token
alone, the compiler can always infer expression types without requiring explicit
type declarations from the programmer. This design ensures that operator
meaning is determined uniquely from the operator token itself, eliminating
semantic ambiguity during type inference and simplifying static analysis.

\subsection{Arithmetic Expressions}
The integer operators \texttt{+}, \texttt{-}, \texttt{*}, and \texttt{/}
require both operands to be \texttt{Integer} and always return \texttt{Integer}
(including division, which uses integer floor division). The float operators
\texttt{+.}, \texttt{-.}, \texttt{*.}, and \texttt{/.} require both operands to
be \texttt{Float} and always return \texttt{Float}. Mixing an \texttt{Integer}
and a \texttt{Float} with any operator is a type error.

\begin{table}[H]
\centering
\caption{Arithmetic operator type rules.\label{tab:arithmetic-result-types}}
\small
\begin{tabular}{llll}
\toprule
Left operand & Operator & Right operand & Result type \\
\midrule
\texttt{Integer} & \texttt{+}, \texttt{-}, \texttt{*}, \texttt{/} & \texttt{Integer} & \texttt{Integer} \\
\texttt{Float}   & \texttt{+.}, \texttt{-.}, \texttt{*.}, \texttt{/.} & \texttt{Float} & \texttt{Float} \\
\texttt{Integer} & any & \texttt{Float} & \texttt{TypeError} \\
\texttt{Float} & any & \texttt{Integer} & \texttt{TypeError} \\
Any & any & \texttt{String} or \texttt{Boolean} & \texttt{TypeError} \\
\bottomrule
\end{tabular}
\end{table}

The following grouped derivation rules capture these constraints:
\[
\frac{\Gamma \vdash e_1 : \texttt{Integer} \qquad \Gamma \vdash e_2 : \texttt{Integer}}
     {\Gamma \vdash e_1\ \mathtt{op}\ e_2 : \texttt{Integer}}
\ \textsc{(T-IntOp)},\ \mathtt{op} \in \{+,\,-,\,*,\,/\}
\]
\[
\frac{\Gamma \vdash e_1 : \texttt{Float} \qquad \Gamma \vdash e_2 : \texttt{Float}}
     {\Gamma \vdash e_1\ \mathtt{op}\ e_2 : \texttt{Float}}
\ \textsc{(T-FloatOp)},\ \mathtt{op} \in \{+.,\,-.,\,*.,\,/.\}
\]
\[
\frac{\Gamma \vdash e_1 : \texttt{Integer} \qquad \Gamma \vdash e_2 : \texttt{Float}}
     {\Gamma \vdash e_1\ \mathtt{op}\ e_2 : \texttt{TypeError}}
\ \textsc{(T-MixedError)}
\]
For the concrete mixed arithmetic expression discussed in this report:
\[
\Gamma \vdash 3 + 4.0 : \texttt{TypeError}
\]
Note that each rule is keyed on the operator token, not on the operand types.
This is only possible because integer and float arithmetic use distinct tokens
--- the property that eliminates operator overloading.

\subsubsection*{Unary negation}
Unary integer negation (\texttt{-x}) is represented internally as the
equivalent subtraction form \texttt{0 - x}. Unary float negation (\verb|-. x|)
is represented internally as \verb|0.0 -. x|. The type checker therefore
applies the standard binary rules
in both cases.

\subsection{Boolean Expressions}
The comparison operators \texttt{==} and \texttt{!=} always produce
\texttt{Boolean}. Both operators require two operands of the \emph{same}
arithmetic type --- either both \texttt{Integer} or both \texttt{Float};
mixed-type comparisons such as \texttt{Integer == Float} are rejected.
Comparisons involving \texttt{Boolean} or \texttt{String} operands are also
rejected by the type checker.
Although \texttt{Boolean} and \texttt{String} values exist in the language,
equality and inequality are deliberately restricted to arithmetic expressions
because the project specification defines basic Boolean expressions as equality
and inequality between two arithmetic expressions.

Furthermore, the same-type requirement for comparisons --- both operands must be
either both \texttt{Integer} or both \texttt{Float} --- is a direct consequence
of the no-overloading design. If \texttt{==} accepted mixed \texttt{Integer} and
\texttt{Float} operands, the compiler would need to determine whether to compare
by promotion (implicitly converting \texttt{Integer} to \texttt{Float}) or by
value identity, introducing the same ambiguity that the separate-operator design
eliminates for arithmetic. Rejecting mixed-type comparisons ensures that
\texttt{==} and \texttt{!=} each have exactly two valid type signatures ---
$(\texttt{Integer}, \texttt{Integer}) \rightarrow \texttt{Boolean}$ and
$(\texttt{Float}, \texttt{Float}) \rightarrow \texttt{Boolean}$ --- rather than a
third mixed-type signature.

\subsection{Assignment Statement}
Assignments define variable types on first use. If a variable does not yet
exist in the type environment, the checker infers the type of the right-hand
side and inserts the variable into the environment. If the variable already
exists, the new value must have the same type.

\[
  \frac{x \notin \Gamma \qquad \Gamma \vdash e : T}
       {\Gamma \vdash x = e : (\Gamma, x:T)}
\]
\[
  \frac{x:T \in \Gamma \qquad \Gamma \vdash e : T}
       {\Gamma \vdash x = e : \Gamma}
\]
\[
  \frac{x:T_1 \in \Gamma \qquad \Gamma \vdash e : T_2 \qquad T_1 \neq T_2}
  {\Gamma \vdash x = e : \texttt{TypeError}}
\]

\subsection{If-Then-Else}
The condition of an \texttt{if} statement must be \texttt{Boolean}. The
then-block and else-block are checked in child scopes so that local operations
can be validated cleanly.

\[
  \frac{\Gamma \vdash c : \texttt{Boolean} \qquad
        \Gamma' \vdash B_{then} \qquad
        \Gamma'' \vdash B_{else}}
       {\Gamma \vdash \texttt{if}(c)\ B_{then}\ \texttt{else}\ B_{else}}
\]
\[
  \frac{\Gamma \vdash c : T \qquad T \neq \texttt{Boolean}}
  {\Gamma \vdash \texttt{if}(c)\ B_{then}\ \texttt{else}\ B_{else} : \texttt{TypeError}}
\]

\subsection{While-Loop}
The condition of a \texttt{while} loop must also be \texttt{Boolean}. The loop
body is checked in a child scope and must not violate the types of variables
already introduced in the surrounding environment.

\[
  \frac{\Gamma \vdash c : \texttt{Boolean} \qquad \Gamma' \vdash B}
       {\Gamma \vdash \texttt{while}(c)\ B}
\]
\[
  \frac{\Gamma \vdash c : T \qquad T \neq \texttt{Boolean}}
  {\Gamma \vdash \texttt{while}(c)\ B : \texttt{TypeError}}
\]

\subsection{Function Definition and Call}
Function definitions are pre-registered in the global scope so that forward
references can be resolved. For functions that are called, parameter types are
fixed at the first call site encountered and then held constant for all
subsequent calls, which must supply arguments of the same types. The return
type is inferred from the \texttt{return} statements in the function body,
which is verified with those fixed parameter types. Inconsistent return types
are reported as a type error.
Furthermore, if a function's inferred return type is non-\texttt{Void}, the
type checker enforces that the function returns on every execution path; a body
where \texttt{return} appears only inside one branch of an \texttt{if} without
a corresponding \texttt{else} is rejected as a non-total function.

For functions that are defined but never called, the type checker performs a
post-pass after all top-level statements have been processed. It infers
parameter types by scanning the body for arithmetic and comparison usage, then
checks the full body with those inferred types. This ensures that type errors
inside never-called function bodies are still detected before execution.

Recursive calls are intentionally outside the scope of this implementation;
supporting them would require additional fixed-point inference or explicit
recursion handling that is orthogonal to the core type-inference algorithm.

\subsubsection*{Formal rules for function inference}
Let $\Sigma$ denote the function-signature environment, where
$\Sigma(f) = (T_1,\ldots,T_n) \rightarrow T_r$.

First-call parameter binding for called functions:
\[
  \frac{\Gamma \vdash a_1 : T_1 \quad \cdots \quad \Gamma \vdash a_n : T_n \qquad f \notin \Sigma}
       {\Sigma \cup \{f:(T_1,\ldots,T_n)\rightarrow \texttt{Void}\},\ \Gamma \vdash f(a_1,\ldots,a_n)}
\ \textsc{(F-FirstCall)}
\]

Subsequent call checking:
\[
  \frac{\Sigma(f)=(T_1,\ldots,T_n)\rightarrow T_r \qquad
        \Gamma \vdash a_1:T_1 \quad \cdots \quad \Gamma \vdash a_n:T_n}
       {\Sigma,\Gamma \vdash f(a_1,\ldots,a_n):T_r}
\ \textsc{(F-Call)}
\]

Return-type consistency:
\[
  \frac{\Sigma(f)=(T_1,\ldots,T_n)\rightarrow \texttt{Void} \qquad
        \Gamma_f \vdash r_1:T \quad \cdots \quad \Gamma_f \vdash r_k:T}
       {\Sigma(f)=(T_1,\ldots,T_n)\rightarrow T}
\ \textsc{(F-ReturnConsistent)}
\]
\[
  \frac{\Gamma_f \vdash r_i:T_i \qquad \Gamma_f \vdash r_j:T_j \qquad T_i \neq T_j}
  {\Gamma_f \vdash \texttt{FunctionDef}(f) : \texttt{TypeError}}
\ \textsc{(F-ReturnMismatch)}
\]

Never-called function fallback used in this implementation:
\[
\frac{p \in \texttt{params}(f) \qquad p\ \text{has no arithmetic-context constraint}}
     {\Sigma(f).p = \texttt{Integer}}
\ \textsc{(F-UncalledDefaultInt)}
\]

Worked derivation (single call):
\begin{lstlisting}
def add(a, b) {
    return a + b;
}

z = add(3, 4);
\end{lstlisting}

\[
\Gamma \vdash 3 : \texttt{Integer}
\qquad
\Gamma \vdash 4 : \texttt{Integer}
\]
\[
\Sigma(add) = (\texttt{Integer},\texttt{Integer}) \rightarrow \texttt{Integer}
\]
\[
\frac{\Gamma \vdash 3 : \texttt{Integer} \qquad \Gamma \vdash 4 : \texttt{Integer}}
  {\Sigma,\Gamma \vdash add(3,4) : \texttt{Integer}}
\]
Hence $\Gamma \vdash z = add(3,4)$ extends $\Gamma$ with
$z:\texttt{Integer}$.

\subsubsection*{Value parameter passing example}
The language uses value parameter passing, or call-by-value: each argument is
evaluated in the caller's environment and a copy of its value is bound to the
corresponding parameter in the callee's environment. Figure~\ref{fig:value-param-passing}
shows the runtime environment before, during, and after a call to a function
that mutates its parameter.

Using semantic-configuration notation with environment $\rho$ and store $\sigma$,
call-by-value is expressed as:
\[
\frac{\begin{aligned}
  &\rho,\sigma \vdash e_i \Downarrow v_i,\sigma_i \quad (1 \le i \le n)\\
  &\rho_f = [x_1 \mapsto v_1,\ldots,x_n \mapsto v_n]\\
  &\rho_f,\sigma_n \vdash B_f \Downarrow \sigma'
\end{aligned}}
{\rho,\sigma \vdash f(e_1,\ldots,e_n) \Downarrow v_{ret},\sigma'}
\qquad \textsc{(S-CallByValue)}
\]
The key point is that each formal parameter $x_i$ is bound to the evaluated
argument value $v_i$ in a fresh activation record for the callee.

The same behavior can be visualized as environment frames:

\begin{table}[H]
\centering
\caption{Activation-record view for one call to \texttt{f}.\label{tab:activation-record-f}}
\small
\begin{tabular}{p{3.4cm}p{8.9cm}}
\toprule
Record field & Content \\
\midrule
Function name & \texttt{f} \\
Caller link & Global frame \\
Parameters & \texttt{a = 5} (copied from caller argument \texttt{x}) \\
Locals after body step & \texttt{a = 6} \\
Return value & \texttt{Void} \\
\bottomrule
\end{tabular}
\end{table}

\begin{table}[H]
\centering
\caption{Runtime environment frames for call-by-value.\label{tab:runtime-frames-cbv}}
\small
\begin{tabular}{p{3.4cm}p{8.9cm}}
\toprule
Phase & Frame contents \\
\midrule
Before call & Global frame: $\{x:5\}$ \\
Call entry & Global frame: $\{x:5\}$; new frame for \texttt{f}: $\{a:5\}$ \\
Inside body & Frame \texttt{f}: assignment \texttt{a = a + 1} updates local frame to $\{a:6\}$ \\
After return & Function frame destroyed; global frame remains $\{x:5\}$ \\
\bottomrule
\end{tabular}
\end{table}

Source program:

\begin{lstlisting}
x = 5;

def f(a) {
    a = a + 1;
}

f(x);
\end{lstlisting}

\begin{figure}[H]
\centering
\begin{tcolorbox}[colback=codebg, colframe=codeframe, boxrule=0.4pt,
                  left=6pt, right=6pt, top=4pt, bottom=4pt]
\footnotesize\ttfamily
\textbf{Before the call}\\[2pt]
Global scope: \quad x $\rightarrow$ 5

\vspace{0.5em}
\textbf{During the call}\\[2pt]
Global scope:\quad x $\rightarrow$ 5\\
Function scope (f): \quad a $\rightarrow$ 5,\ then a $\rightarrow$ 6

\vspace{0.5em}
\textbf{After the call}\\[2pt]
Global scope: \quad x $\rightarrow$ 5
\end{tcolorbox}
\caption{Value parameter passing: the parameter \texttt{a} receives a copy of
\texttt{x}, so mutating \texttt{a} inside \texttt{f} does not affect
\texttt{x} in the caller.\label{fig:value-param-passing}}
\end{figure}

A parameter receives a copy of the argument value rather than a
shared reference, modifications to \texttt{a} inside the function do not
affect \texttt{x} in the caller environment. This demonstrates standard
call-by-value semantics and avoids aliasing effects associated with
reference-based mechanisms.

\section{Type Error Reporting}
Type errors are reported with source positions so that the user can locate the
problem quickly. The format used by the system is:

\begin{center}
\texttt{[line X, col Y] TypeError: message}
\end{center}

Examples include assigning a \texttt{String} to an \texttt{Integer} variable,
using a non-Boolean condition in an \texttt{if} statement, or supplying the
wrong argument type to a function.

\begin{table}[H]
\centering
\caption{Representative type errors detected by the static type checker.\label{tab:type-errors}}
\small
\begin{tabular}{p{5.4cm}p{5.4cm}p{2.4cm}}
\toprule
Invalid program & Reason & Detection stage \\
\midrule
\verb|x = 5; x = "hello";| & Inconsistent variable type & Type checker \\
\texttt{3 + 4.0}             & Mixed arithmetic types     & Type checker \\
\texttt{if (3) \{ \ldots \}} & Condition not Boolean      & Type checker \\
\verb|print()|               & Wrong argument count       & Type checker \\
\verb|x = print(5);|         & Void assigned to variable  & Type checker \\
\texttt{true == false}       & Boolean comparison unsupported & Type checker \\
\bottomrule
\end{tabular}
\end{table}

\section{Example Type Inference Traces}
Consider the following program fragment:

\begin{lstlisting}
x = 3;
y = 4;
z = x + y;
print(z);
\end{lstlisting}

The checker infers:
\begin{itemize}
	\item \texttt{x: Integer} (from literal \texttt{3})
	\item \texttt{y: Integer} (from literal \texttt{4})
	\item \texttt{x + y: Integer} (Integer \texttt{+} Integer $\rightarrow$ Integer)
	\item \texttt{z: Integer}
\end{itemize}

For a float example using dot-suffixed operators:

\begin{lstlisting}
def addf(a, b) {
		return a +. b;
}

result = addf(2.0, 3.5);
\end{lstlisting}

The first function call infers \texttt{a: Float} and \texttt{b: Float}. The
operator \texttt{+.} requires both operands to be \texttt{Float}, so the
function return type is inferred as \texttt{Float}. The variable
\texttt{result} is also inferred as \texttt{Float}.

\section{Static Semantics Summary}

Table~\ref{tab:static-semantics} summarises the key typing rules enforced by
the type checker.

\begin{table}[H]
\centering
\caption{Static semantics summary: constructs and their typing rules.\label{tab:static-semantics}}
\small
\begin{tabular}{p{5.5cm}p{7.5cm}}
\toprule
Construct & Rule \\
\midrule
Integer arithmetic ($+$, $-$, $*$, $/$) &
  \texttt{Integer} $\times$ \texttt{Integer} $\rightarrow$ \texttt{Integer} \\
Float arithmetic (\texttt{+.}, \texttt{-.}, \texttt{*.}, \texttt{/.}) &
  \texttt{Float} $\times$ \texttt{Float} $\rightarrow$ \texttt{Float} \\
Equality / inequality (\texttt{==}, \texttt{!=}) &
  Same numeric type $\rightarrow$ \texttt{Boolean}; Boolean/String operands rejected \\
Assignment &
  RHS type must match the variable's inferred type \\
\texttt{if} / \texttt{while} condition &
  Condition expression must be \texttt{Boolean} \\
Function return &
  All \texttt{return} expressions must share one type; that type becomes the function type \\
\bottomrule
\end{tabular}
\end{table}
